\lecture{8.11}{Dynamic Treatment Rules}
\label{lec:8-11}

\subsection{Static vs Dynamic Regimes}

A static regime assigns the same treatment sequence to everyone:
\[
  \abar=(a_0,\ldots,a_K).
\]
Examples are always treat, never treat, or treat at visits 0 and 1 only.

A dynamic regime assigns treatment according to observed history. It is a
sequence of rules
\[
  d=(d_0,\ldots,d_K),
\]
where
\[
  A_k=d_k(\Lbar_k,\Abar_{k-1}).
\]
The rule may depend on laboratory values, symptoms, adverse events, previous
treatment, or other information available at time $k$.

\subsection{Rule-Based Interventions}

\begin{definition}[Dynamic treatment rule]
A dynamic treatment rule is an intervention that sets treatment at each time to
a deterministic or stochastic function of the individual's observed history up
to that time.
\end{definition}

For example, in inflammatory bowel disease, a drug may be recommended only while
liver function remains adequate:
\[
  d_k(\Lbar_k)=
  \begin{cases}
  1, & \text{if recent liver function is adequate},\\
  0, & \text{otherwise}.
  \end{cases}
\]
This rule is more realistic than ``always treat'' if treatment is contraindicated
after signs of liver damage.

\subsection{Treatment Decisions Based on Evolving Covariates}

Dynamic regimes are often clinically natural because real treatment decisions
respond to evolving information. A rule might say:
\begin{itemize}
  \item continue chemotherapy if tumor response is adequate and toxicity is mild;
  \item intensify antihypertensive therapy if blood pressure remains high;
  \item stop a medication if a severe adverse event occurs;
  \item initiate rescue therapy if symptoms cross a threshold.
\end{itemize}

\begin{figure}[H]
\centering
\begin{tikzpicture}[
  node distance=1.05cm,
  box/.style={draw, rounded corners=2pt, minimum width=1.15cm, minimum height=.72cm},
  rule/.style={draw, rounded corners=2pt, fill=gray!8, minimum width=2.1cm, minimum height=.72cm},
  arr/.style={-{Latex[length=2.2mm]}, thick}
]
\node[box] (L0) {$L_0$};
\node[rule, right=of L0] (d0) {$A_0=d_0(L_0)$};
\node[box, right=of d0] (L1) {$L_1$};
\node[rule, right=of L1] (d1) {$A_1=d_1(\Lbar_1,A_0)$};
\node[right=of d1] (dots) {$\cdots$};
\node[box, right=of dots] (Y) {$Y(d)$};
\draw[arr] (L0) -- (d0);
\draw[arr] (d0) -- (L1);
\draw[arr] (L1) -- (d1);
\draw[arr] (d1) -- (dots);
\draw[arr] (dots) -- (Y);
\end{tikzpicture}
\caption{A dynamic rule maps the evolving observed history into treatment decisions.}
\label{fig:ch8-dynamic-rule}
\end{figure}

\subsection{Counterfactual Outcomes Under Dynamic Rules}

Let $Y(d)$ denote the outcome that would be observed if treatment were assigned
according to rule $d$ at every time point. The causal mean under the rule is
\[
  \psi(d)=\E\{Y(d)\}.
\]
A causal contrast between two rules is
\[
  \psi(d)-\psi(d')
  =
  \E\{Y(d)\}-\E\{Y(d')\}.
\]

The consistency assumption for a dynamic rule states that if an individual's
observed treatment history happens to agree with the rule at every time, then
the observed outcome equals the counterfactual outcome under that rule:
\[
  \text{if } A_k=d_k(\Lbar_k,\Abar_{k-1})\text{ for all }k,
  \quad Y=Y(d).
\]

\subsection{Identification Extensions}

Identification conditions are the dynamic-rule versions of the static-regime
conditions. Sequential exchangeability becomes
\[
  Y(d) \ci A_k \given \Lbar_k,\Abar_{k-1}
\]
among histories compatible with the rule up to time $k-1$. Positivity requires
positive probability of receiving the rule-recommended treatment:
\[
  \Prob\{A_k=d_k(\Lbar_k,\Abar_{k-1})
  \given \Lbar_k,\Abar_{k-1}\}>0
\]
for histories that can occur under the rule.

The g-computation formula is obtained by replacing fixed values $a_k$ with rule
values $d_k(\Lbar_k,\Abar_{k-1})$. For example, the time-specific treatment
mechanism used in IPTW is evaluated at the rule-recommended treatment:
\[
  g_k^d(\Lbar_k,\Abar_{k-1})
  =
  \Prob\{A_k=d_k(\Lbar_k,\Abar_{k-1})
  \given \Lbar_k,\Abar_{k-1}\}.
\]
The dynamic-regime weight is
\[
  H_d=
  \prod_{k=0}^K
  \frac{\ind\{A_k=d_k(\Lbar_k,\Abar_{k-1})\}}
  {g_k^d(\Lbar_k,\Abar_{k-1})}.
\]

\begin{identificationbox}
Under consistency, sequential exchangeability, and positivity for the dynamic
rule $d$,
\[
  \E\{Y(d)\}
  =
  \E\left[
    \prod_{k=0}^K
    \frac{\ind\{A_k=d_k(\Lbar_k,\Abar_{k-1})\}}
    {\Prob\{A_k=d_k(\Lbar_k,\Abar_{k-1})
      \given \Lbar_k,\Abar_{k-1}\}}
    Y
  \right],
\]
and an analogous dynamic g-computation formula is obtained by recursively
setting each treatment to the value recommended by $d$.
\end{identificationbox}

Dynamic regimes can reduce positivity problems because they avoid treatment
assignments that are impossible or clinically inappropriate for some histories.
For example, ``treat while liver function permits'' may be identifiable in a
population where ``always treat regardless of liver damage'' is not.

\subsection{Optimal Dynamic Regimes}

A separate line of work asks which rule among a class of candidate rules
maximizes a desired outcome:
\[
  d_{\mathrm{opt}}\in \argmax_{d\in\mathcal{D}}\E\{Y(d)\}.
\]
This is an optimal dynamic treatment regime problem. It is harder than
estimating the value of a pre-specified rule because the same data are used to
select the rule and estimate its value. Valid inference must account for this
selection step.

\begin{keyidea}
Dynamic treatment rules replace fixed treatment values with history-dependent
decisions. They are often more realistic, more clinically interpretable, and
less vulnerable to structural positivity violations than static regimes.
\end{keyidea}
